\documentclass[10pt,a4paper]{amsart}
\usepackage[margin=19mm]{geometry}
\usepackage{amsmath,amssymb,mathtools}
\usepackage[scr=rsfs,cal=euler]{mathalfa}
\usepackage{xfrac}
\usepackage[unicode,hidelinks,bookmarks=false]{hyperref}
\newcommand{\ie}{i.e.\ }
\newcommand{\cf}{cf.\ }
\newcommand{\resp}{respectively\ }
\newcommand{\Subsection}{\subsection}
\providecommand{\nobreakdash}{\nobreak\hbox{-}}
\setlength{\emergencystretch}{2em}
\multlinegap0pt
\usepackage{bm}
\usepackage{bbm}
\usepackage{dsfont}
\usepackage{xspace}
\usepackage{cancel}
\usepackage{mathtools}
\usepackage{amsthm}
\makeatletter
\@ifundefined{theorem}{\newtheorem{theorem}{Theorem}}{}
\@ifundefined{lemma}{\newtheorem{lemma}[theorem]{Lemma}}{}
\makeatother
\newcommand{\D}{\mathbb{D}}
\newcommand{\SSS}{\mathbb{S}}
\newcommand{\T}{\mathbb{T}}
\newcommand{\Z}{\mathbb{Z}}
\title[Index theory and bulk-boundary correspondence for inversion-]{Index theory and bulk-boundary correspondence for inversion-symmetric second-order topological insulators}
\author{Shin Hayashi}
\date{Source: arXiv:2509.09240v1 (CC BY 4.0)}
\begin{document}
\maketitle

\noindent\emph{Source and licence.} Index theory and bulk-boundary correspondence for inversion-symmetric second-order topological insulators. Version arXiv:2509.09240v1 (CC BY 4.0), distributed under the Creative Commons Attribution 4.0 International licence, as recorded in the arXiv licence metadata for this exact version and shown on the abstract page https://arxiv.org/abs/2509.09240v1. Full rights notice: licenses/arxiv-2509.09240-CC-BY-4.0.txt.\par\smallskip

\noindent\textit{Editorial scope.} Counted unit: the Lemma whose source label is \texttt{mape} in the article (source0.tex lines 1198--1201, source label \texttt{mape}), delivered together with the article's own complete original proof of it (source0.tex lines 1203--1221, one \texttt{proof} environment of 19 lines). No mathematics is written, repaired or re-derived: statement and proof are the article's own text line for line. Required context retained uncounted: X (lines 615--621); Spm (lines 643--646); KX (lines 1166--1169). Every definition, display and label of those lines is retained unchanged. Omitted from this package and not redistributed: 1--614, 622--642, 647--1165, 1170--1197, 1202--1202, 1222--2501. Nothing inside a retained excerpt is omitted or altered: every retained body is delivered exactly as the article's lines are written, so no omission receipt of any kind accompanies this package. Citations: the retained excerpts carry no citation command at all, so no bibliography entry of the article is transcribed and no reference is redistributed. Reference adaptation: none was needed and none was made; every reference inside the delivered unit resolves inside it, so no unresolved reference or citation is produced. Printed slips kept literal: no printed word, symbol, hypothesis or number of the counted statement or of its proof was altered. Layout and class: the article's own class is \texttt{amsart} with packages including amssymb, amsmath, amsthm, amscd, amstext, amsfonts, xy, graphicx; the wrapper is the lane's pinned \texttt{amsart} editorial wrapper at 10pt on A4, which restates the theorem-like environments the retained text needs and adds the article's own macro definitions that the retained text needs (\texttt{\textbackslash D}, \texttt{\textbackslash SSS}, \texttt{\textbackslash T}, \texttt{\textbackslash Z}), with extra wrapper packages (bm, bbm, dsfont, xspace, cancel, mathtools). The delivered numbering is the unit's own. Assets: no retained span contains a figure environment, a graphics-inclusion command, a PDF-page inclusion, a table or a TikZ picture; no image file is copied into this package, the package declares no asset dependency and no image file is redistributed with it. Licence: version arXiv:2509.09240v1 of this article is distributed under the Creative Commons Attribution 4.0 International licence, as recorded in the arXiv licence metadata for this exact version and shown on the abstract page https://arxiv.org/abs/2509.09240v1. A full rights notice is delivered at licenses/arxiv-2509.09240-CC-BY-4.0.txt. Characters: every retained excerpt is pure ASCII, so the delivered engine, XeLaTeX with its default Latin Modern text font, reports no missing character.
\begin{equation}
\label{X}
	X = (\D_+ \times \T) \underset{\T^2}{\cup}
		(\D_- \times \T) \underset{\T^2}{\cup}
			(\T \times \D_+)	 \underset{\T^2}{\cup}
				(\T \times \D_-).
\end{equation}

\begin{equation}
\label{Spm}
	\SSS_{\pm\pm}^3 = (\D_{\pm} \times \T) \underset{\T^2}{\cup} (\T \times \D_{\pm}),
\end{equation}

\begin{lemma}
\label{KX}
	$K^0(X) \cong \Z^3$, $K^1(X) \cong \Z^3$.
\end{lemma}

\begin{lemma}
\label{mape}
The map $e = \sum_{i=1}^3 e_i \colon \oplus_{i=1}^3 K^1(\overline{B}_i, \partial \overline{B}_i) \to K^1(X)$ is an isomorphism.
\end{lemma}

\begin{proof}
Let $i_{\pm \pm} \colon \SSS_{\pm\pm}^3 \hookrightarrow X$ be the inclusion where double signs correspond (see (\ref{X}) and (\ref{Spm})), and
let us consider the following maps:
\begin{equation}
\oplus_{i=1}^3 K^1(\overline{B}_i, \partial \overline{B}_i) \overset{e}{\to}
	K^1(X) \overset{i}{\to}
		K^1(\SSS^3_{++}) \oplus K^1(\SSS^3_{-+}) \oplus K^1(\SSS^3_{--})
\end{equation}
where $i = i_{++}^* \oplus i_{+-}^* \oplus i_{--}^*$.
By using the orientations of $B_i$ and $\SSS^3_{\pm \pm}$, both groups $\oplus_{i=1}^3 K^1(\overline{B}_i, \partial \overline{B}_i)$ and $K^1(\SSS^3_{++}) \oplus K^1(\SSS^3_{-+}) \oplus K^1(\SSS^3_{--})$ are identified with $\Z^3$ via the three-dimensional winding numbers.
Through these identifications, the map $i \circ e$ is computed as follows\footnote{Note that our open three-ball $B_i$ is contained in some solid torus, and that each of our three-spheres $\SSS^3_{\pm\pm}$ is the union of two solid tori. The orientation of the three-ball $B_i$ can be the same as or opposite to that of the three-sphere $\SSS^3_{\pm\pm}$, which leads to a positive or negative sign in (\ref{sign}), respectively.}:
\begin{equation}
\label{sign}
	i \circ e (p,q,r) = (p+q, -q+r, -r),
\end{equation}
where $p,q,r \in \Z$, which is a bijection.
Since $i$ is surjective and $K^1(X)$ is isomorphic to $\Z^3$ by Lemma~\ref{KX}, the map $i$ is bijective.
Therefore, $e$ is also bijective.
\end{proof}

\end{document}
